\section{Ordinary signature depth, staircase projection, and ridge geometry}
\label{app:depth-and-ridge}

Appending a clock or frequency coordinate to an uncoded path makes positional
information available in the full infinite signature. Two mathematical questions
connect that fact to the depth-two Pascal lift. First, when does a bounded-depth
ordinary signature recover the depth-two Pascal coordinates linearly? Second,
when two signature representations span the exact same affine linear function
space on a dataset, why can a linear ridge classifier still make different
predictions? We answer both questions in turn.

\subsection{Clocked paths at depth three and the interpolation correction}
In Example~\ref{ex:two-rises}, the uncoded clock-augmented paths \((T, x)\) for
the spread schedule \(a = (1, 0, 0, 1)\) and central schedule \(b = (0, 1, 1, 0)\)
had identical coordinates through depth two, while the two-copy Pascal area
\(J_{01}\), between the original and weighted copies with increments
\((z_r,u_r z_r)\), separated them as \(1/2\) versus \(1/6\). Moving one level higher
in the ordinary signature of \((T, x)\) on \([0, 1]\) exposes the space--time--space
iterated integral:
\[
 S^{xTx} = \int_{0 < t_1 < t_2 < t_3 < 1} dx_{t_1}\,dT_{t_2}\,dx_{t_3}
         = \int_0^1 \Bigl(\int_0^{t_2} dx_{t_1}\Bigr)\Bigl(\int_{t_2}^1 dx_{t_3}\Bigr)\,dt_2
         = \int_0^1 x_t (s_x - x_t)\,dt,
\]
since \(dT_t = dt\) and \(x_0 = 0\). With total rise \(s_x = 2\) and four uniform
clock steps of \(1/4\), integrating \(x_t(2 - x_t)\) across the four piecewise-linear
segments gives
\[
 S^{xTx}(a) = \frac{5}{6}, \qquad S^{xTx}(b) = \frac{1}{3},
\]
and the corresponding degree-three logsignature coefficients of the word \(xTx\)
are \(1/6\) for \(a\) and \(-1/3\) for \(b\). The ordinary clocked path therefore
first distinguishes the two schedules at depth three.

To see the exact relationship between \(S^{xTx}\) and the discrete separation
area \(J_{01}\) for an arbitrary scalar increment sequence
\(z=(z_r)_{r=0}^{R-1}\), with \(z_r\in\R\), over \(R \geq 2\) uniform clock steps
\(\Delta T_r = 1/R\), decompose \(S^{xTx}\) over the segment indices
\(r, m, t \in \{0, \ldots, R-1\}\) in which the three integration times
\(t_1 < t_2 < t_3\) fall:
\begin{itemize}
\item \emph{Distinct spatial segments (\(r < t\)):} If \(t_2\) falls in a
strictly intermediate segment \(r < m < t\) (of which there are \(t - r - 1\)),
the integral contributes \(\frac{1}{R} z_r z_t\) per segment. If \(t_2\) falls in
segment \(r\) with \(t_1 < t_2\), the ordered \(2\)-simplex on segment \(r\)
contributes \(\frac{1}{2R} z_r z_t\); symmetrically, \(t_2 < t_3\) in segment
\(t\) contributes another \(\frac{1}{2R} z_r z_t\). Summing these three cases
yields \(\bigl((t - r - 1) + \tfrac12 + \tfrac12\bigr)\frac{1}{R} z_r z_t = \frac{t - r}{R} z_r z_t\).
\item \emph{Within-segment triples (\(r = m = t\)):} When all three integration
variables fall inside the same segment \(r\), linear traversal of
\((\Delta T_r, z_r) = (1/R, z_r)\) over the ordered \(3\)-simplex of volume
\(1/3! = 1/6\) contributes \(\frac{1}{6R} z_r^2\).
\end{itemize}
Summing over all segments and comparing with
\(J_{01} = \frac{1}{2(R-1)}\sum_{r < t}(t - r) z_r z_t\) from~\eqref{eq:app-separation}
establishes the exact identity
\begin{equation}\label{eq:discrete-depth-correction}
 J_{01} = \frac{R}{2(R-1)} S^{xTx}
             - \frac{1}{12(R-1)} \sum_{r=0}^{R-1} z_r^2.
\end{equation}
For both \(a\) and \(b\), \(\sum_{r=0}^3 z_r^2 = 2\), so the within-segment
interpolation correction is \(2/(12 \cdot 3) = 1/18\), confirming
\(\frac{4}{6}\cdot\frac{5}{6} - \frac{1}{18} = \frac{1}{2}\) and
\(\frac{4}{6}\cdot\frac{1}{3} - \frac{1}{18} = \frac{1}{6}\).

\begin{example}[Depth-three collision under simultaneous linear interpolation]
\label{ex:depth-three-collision}
Because the within-step sum of squares \(\sum_{r=0}^{R-1} z_r^2\) is not a
function of the ordinary depth-three signature of \((T, x)\),
identity~\eqref{eq:discrete-depth-correction} prevents any universal recovery of
\(J_{01}\) from clocked depth-three coordinates. Explicitly, on the
four-step uniform clock \(\Delta T_r = 1/4\), consider the two scalar increment
sequences
\[
 z^{(1)} = (2, 1, 1, 1), \qquad
 z^{(2)} = \Bigl(\frac{17}{8},\, \frac{5}{8},\, \frac{11}{8},\, \frac{7}{8}\Bigr).
\]
Direct evaluation of the concatenated segment exponentials shows that the clocked
paths of \(z^{(1)}\) and \(z^{(2)}\) have \emph{identical} ordinary tensor
signatures at every word of length zero, one, two, and three (all \(15\) tensor
coordinates coincide, with \(s = (1, 5)\) and \(S^{xTx} = 103/24\)). Yet their sum
of squared increments differs (\(7\) versus \(121/16\)), so~\eqref{eq:discrete-depth-correction}
assigns them distinct depth-two Pascal areas:
\[
 J_{01}\bigl(z^{(1)}\bigr) = \frac{8}{3} = \frac{512}{192}, \qquad
 J_{01}\bigl(z^{(2)}\bigr) = \frac{509}{192}.
\]
Conversely, discrete index coding depends on how steps are segmented and
anchored: prepending an idle zero step \((0, 0)\) to \(a\) and \(b\) leaves the
geometric curve \((T, x)\) and all of its ordinary signature coordinates
unchanged, while changing the five-step Pascal grid to \(u_r = r/4\) and
rescaling their coded areas from \((1/2, 1/6)\) to \((3/8, 1/8)\).
\end{example}

\subsection{Exact scalar projection from a frequency--mass staircase}
\label{app:staircase-projection}
The discrepancy in~\eqref{eq:discrete-depth-correction} arose because \(T\) and
\(x\) changed simultaneously within each piecewise-linear segment. When a
one-dimensional sequence \(p_0, \ldots, p_{K-1}\) is instead embedded into \(\R^2\)
as a \emph{staircase} that holds the position coordinate constant during each
data jump and advances position only between jumps, within-step position changes
vanish identically, and every scalar Pascal-coded coordinate becomes an exact
linear combination of higher-depth staircase signature coordinates.

Let \(K \geq 3\) and \(u_k = k/(K-1)\) for \(0 \leq k < K\). Construct the
two-channel frequency--mass staircase \((F, M)\) starting at \((0, 0)\) as
follows: at bin \(k = 0\), traverse a vertical step \(\Delta M = p_0\) while
\(F = u_0 = 0\) remains fixed; for each \(k = 0, \ldots, K-2\), move horizontally
by \(\Delta F = u_{k+1} - u_k = 1/(K-1)\) with \(\Delta M = 0\), then traverse
the vertical mass step \(\Delta M = p_{k+1}\) at fixed \(F = u_{k+1}\). No
horizontal step follows the final jump \(p_{K-1}\) at \(u_{K-1} = 1\), so for a
unit-sum spectrum the endpoint is \((F_{\mathrm{end}}, M_{\mathrm{end}}) = (1, 1)\).

In signature words over the alphabet \(\{F, M\}\), let \(F^a\) denote \(a\)
consecutive \(F\) letters (with \(F^0\) the empty word). Write \(s_a\) and
\(B_{ab}\) for the first- and second-level coordinates of the scalar lift with
monomial increment copies \(u_k^a p_k\) and \(u_k^b p_k\), and set
\(J_{ab} = \tfrac12(B_{ab} - B_{ba})\). Here \(a,b\) label polynomial copies,
while \(k\) indexes frequency bins; in particular, \(J_{01}\) pairs the original
copy with the linearly weighted copy.

The following result specializes the standard polynomial one-form representation
of iterated integrals as signature functionals~\cite[Lem.~4.1]{polynomialoneforms},
applied to \(F^a\,dM\). We give explicit coefficients and the sufficient scalar
depth bound \(2P-1\).

\begin{theorem}[Exact scalar staircase projection]\label{prop:staircase-projection}
Along the \((F, M)\) staircase on \(K \geq 3\) bins, the monomial-coded endpoints
and signed areas for degrees \(0, 1, 2\) satisfy
\begin{align}
 s_0 &= S^M, & s_1 &= S^{FM}, & s_2 &= 2 S^{FFM}, \label{eq:staircase-endpoints}\\
 J_{01} &= \frac{1}{2} S^{MFM}, &&& \label{eq:staircase-area01}\\
 J_{02} &= S^{FMFM} + S^{MFFM}, &&& \label{eq:staircase-area02}\\
 J_{12} &= S^{FFMFM} + S^{FMFFM}. &&& \label{eq:staircase-area12}
\end{align}
Consequently, at any fixed \(K \geq 3\), all depth-two endpoints and signed areas
of the \(P = 3\) Pascal lift---as well as its positive-diagonal QR version
(Section~\ref{app:qr}) and \(\sqrt 2\)-scaled spectral areas---are fixed,
data-independent linear projections of the ordinary staircase tensor signature
through depth five. More generally, for \(2 \leq P \leq K\) scalar polynomial
copies of degrees below \(P\), ordinary staircase depth \(2P - 1\) suffices.
\end{theorem}
\begin{proof}
During the \(k\)th vertical segment, \(dM\) integrates to \(p_k\) at constant
frequency \(F_t = u_k\), whereas \(dM = 0\) on every horizontal segment. Since
\(F_0 = 0\), the ordered \(a\)-fold integral of \(dF\) from \(0\) up to time
\(t\) equals \(F_t^a / a!\). Integrating against \(dM_t\) therefore yields
\[
 a!\,S^{F^a M} = \int_0^T F_t^a\,dM_t = \sum_{k=0}^{K-1} u_k^a p_k = s_a,
\]
which gives~\eqref{eq:staircase-endpoints} for \(a = 0, 1, 2\).

For the second-level coordinate \(B_{ab} = \int_{0 < s < t < T} F_s^a F_t^b\,dM_s\,dM_t\)
(notice that because \(F\) is constant within each vertical mass step, integrating
\(F_s^a F_t^b\,dM_s\,dM_t\) over \(s < t\) inside the \(k\)th vertical segment
produces \(\tfrac12 u_k^{a+b} p_k^2\), matching the piecewise-linear diagonal
term in~\eqref{eq:app-depth-two}), expand \(F_t^b = \bigl(F_s + (F_t - F_s)\bigr)^b\)
by the binomial theorem:
\[
 F_s^a F_t^b = \sum_{j=0}^b \binom{b}{j} F_s^{a+b-j} (F_t - F_s)^j.
\]
For \(0 < s < t\), \(F_s^{a+b-j} / (a+b-j)!\) is the \((a+b-j)\)-fold iterated
integral of \(dF\) over \((0, s)\), and \((F_t - F_s)^j / j!\) is the \(j\)-fold
iterated integral of \(dF\) over \((s, t)\). Interleaving these \(dF\) factors
with \(dM_s\) and \(dM_t\) yields, for all integers \(a, b \geq 0\),
\begin{equation}\label{eq:staircase-general}
 B_{ab} = \sum_{j=0}^b \binom{b}{j} (a + b - j)!\,j!\,S^{F^{a+b-j} M F^j M}.
\end{equation}
Notice that the \(j = 0\) term \((a+b)!\,S^{F^{a+b} M M}\) is symmetric in \(a\)
and \(b\) and cancels upon forming \(J_{ab} = \tfrac12(B_{ab} - B_{ba})\).
Evaluating~\eqref{eq:staircase-general} at \((a, b) = (0, 1), (0, 2), (1, 2)\)
and subtracting \(B_{ba}\) gives:
\begin{align*}
 B_{01} - B_{10} &= (S^{FMM} + S^{MFM}) - S^{FMM} = S^{MFM}, \\
 B_{02} - B_{20} &= (2S^{FFMM} + 2S^{FMFM} + 2S^{MFFM}) - 2S^{FFMM}
                  = 2(S^{FMFM} + S^{MFFM}), \\
 B_{12} - B_{21} &= (6S^{FFFMM} + 4S^{FFMFM} + 2S^{FMFFM})
                  - (6S^{FFFMM} + 2S^{FFMFM}) \\
                 &= 2(S^{FFMFM} + S^{FMFFM}),
\end{align*}
proving~\eqref{eq:staircase-area01}--\eqref{eq:staircase-area12}.

Finally, any degree-\((P-1)\) polynomial basis \(c_p(u) = \sum_{a=0}^{P-1} L_{pa} u^a\)
transforms endpoints and areas linearly via \(s^c = L s\) and \(J^c = L J L^\top\)
by~\eqref{eq:basis-naturality}, including the normalized Pascal weights, the
orthonormal QR weights \(q(k) = (D H^{-1})^\top b(k)\) of Section~\ref{app:qr},
and multiplication of areas by \(\sqrt 2\). For \(0 \leq a < b \leq P - 1\), the
words in~\eqref{eq:staircase-general} have length \((a + b - j) + 1 + j + 1 = a + b + 2 \leq (P - 2) + (P - 1) + 2 = 2P - 1\),
while endpoint words \(F^a M\) have length \(a + 1 \leq P < 2P - 1\) for
\(P \geq 2\).
\end{proof}

\begin{remark}[Scope of the staircase construction]
Every word appearing in~\eqref{eq:staircase-endpoints}--\eqref{eq:staircase-general}
ends in the letter \(M\). Appending an extra horizontal step after the last mass
jump would change \(S^F\) and words ending in \(F\), while leaving every formula
in Theorem~\ref{prop:staircase-projection} intact. By contrast, replacing the
axis-aligned staircase with diagonal segments that increment \(F\) and \(M\)
simultaneously reintroduces within-bin interpolation terms analogous
to~\eqref{eq:discrete-depth-correction}. Likewise, Theorem~\ref{prop:staircase-projection}
exploits the scalar structure of \(M\) (where same-degree areas \(J_{aa}\) vanish)
and does not assert an identical depth-\((2P-1)\) projection for general
multichannel temporal paths.
\end{remark}

\subsection{Fixed endpoints, affine spans, and ridge regularization}
\label{app:fixed-endpoints}
Theorem~\ref{prop:staircase-projection} shows that the five nonconstant \(P=3\)
spectral Pascal coordinates lie inside the linear span of the staircase tensor
signature through depth five. When comparing candidate representations empirically
with a regularized linear classifier, however, matching or enlarging the linear
span is not equivalent to preserving ridge predictions.

The standard tensor exponential and logarithm identities~\cite{chen} show that,
even at depth three, fixed endpoints create exact affine equivalences between
tensor and logsignature coordinates. Writing \(\ell_1 = s\), \(\ell_2 = J\), and
\(\ell_3\) for the homogeneous components of the tensor logarithm
\(\log_\otimes S_{\leq 3}(y)\) through degree three, expanding the tensor
exponential \(S_{\leq 3}(y) = \exp_\otimes(\ell_1 + \ell_2 + \ell_3)\) gives
\begin{equation}\label{eq:fixed-endpoint-span}
 S^2 = J + \frac{1}{2} s^{\otimes 2}, \qquad
 S^3 = \ell_3 + \frac{1}{2}(s \otimes J + J \otimes s) + \frac{1}{6} s^{\otimes 3}.
\end{equation}
On a family of paths whose endpoint \(s\) is identical for every observation---such
as the unit-sum normalized \((F, M)\) staircases of Subsection~\ref{app:staircase-projection},
all of which share \(s = (1, 1)\)---the relations~\eqref{eq:fixed-endpoint-span}
and their Lie-series inverses are \emph{affine linear} in the remaining
coordinates \((J, \ell_3)\) and \((S^2, S^3)\). Consequently, through depth three,
staircase tensor features and staircase logsignature features span the exact same
affine linear function space on unit-sum spectra. (At depth four, the exponential
expansion of \(S^4\) introduces the quadratic area product \(\tfrac12 J \otimes J\),
so fixed endpoints alone no longer guarantee affine equivalence with the
depth-four logsignature.)

Why, then, can two feature maps with the same affine span produce different
classifications under a linear ridge model? Suppose \(Z \in \R^{N \times d_1}\)
and \(\widetilde Z \in \R^{N \times d_2}\) are two column-standardized training
feature matrices having the exact same column space \(\operatorname{col}(Z) = \operatorname{col}(\widetilde Z) \subseteq \{u \in \R^N : \mathbf 1^\top u = 0\}\).
For centered targets \(\widetilde Y = Y - \mathbf 1 \bar y^\top\) and ridge
penalty \(\alpha > 0\), solving
\[
 \min_{W} \|\widetilde Y - Z W\|_F^2 + \alpha \|W\|_F^2
\]
yields the standard dual ridge predictions~\cite{ridge}
\(\widehat Y = Z(Z^\top Z + \alpha I_{d_1})^{-1} Z^\top \widetilde Y = K(K + \alpha I_N)^{-1} \widetilde Y\),
determined entirely by the Gram matrix \(K = Z Z^\top\). At fixed \(\alpha\),
the two training prediction operators agree for all centered targets exactly
when \(K = \widetilde K = \widetilde Z \widetilde Z^\top\). For test features
\(Z_*\) standardized using the training statistics, the scores are
\(Z_* Z^\top (K + \alpha I_N)^{-1}\widetilde Y + \mathbf 1\bar y^\top\).
Columnwise standardization gives each nonconstant training column norm
\(\sqrt N\), while the Gram matrices can still differ.
\begin{enumerate}
\item \emph{Redundant or linearly dependent coordinates:} If \(\widetilde Z\)
repeats a standardized column \(k\) times (for example, when \(S^{FM}\) and
\(S^{MF} = 1 - S^{FM}\) both standardize to \(\pm z_1\) on unit-endpoint
staircases), that direction receives \(k\) times larger weight in
\(\widetilde Z \widetilde Z^\top\), effectively dividing its ridge penalty by
\(k\) relative to unrepeated directions.
\item \emph{Non-orthogonal basis transforms and unprojected ambient dimensions:}
Replacing \(P=3\) QR-Pascal coordinates by unorthogonalized monomials can alter
the singular values of \(Z\) without changing its column space. Fitting ridge
directly to all \(1{,}040\) depth-five staircase tensor columns can also add
directions beyond the \(100\) projected Pascal features under the same isotropic
penalty \(\alpha \|W\|_F^2\).
\end{enumerate}
Part~II's empirical comparisons---especially the staircase depth table in
Section~\ref{app:har} and the width-varying interaction controls in
Section~\ref{app:natops}---directly reflect this interplay between expressive span
and ridge metric geometry.
